\textsc{\cref{mainnonsenseprojective,mainnonsenseinjective,projcorrespondence} give general reconstruction results}
in the non-rigid setting,
which has recently seen increased interest in multiple areas,
see for example~\cite{douglas19,allen21:dualit-feigin-fuchs} or~\cite[Corollary~10.11]{stroinski24:modul-tambar}.
However, lax \(\cat{C}\)\hyp{}module monads on \(\cat{C}\) are less accessible than mere algebra objects of \(\cat{C}\),
and may indeed appear not very accessible in general.
To show that this need not be the case, we again turn to the Hopf-algebraic setting.

Recall that a ring category%
—the non-rigid counterpart of a tensor category, see~\cite[Sections~4.2 and~5.4]{Etingof2015}—%
which admits a fibre functor to \(\tetramodfd\) is monoidally equivalent to the category
of finite-dimensional left \(B\)\hyp{}comodules over a (not necessarily Hopf) bialgebra \(B\).
In this chapter, we show that the category of left exact finitary lax \(\Tetramod[B][][][]\)\hyp{}module endofunctors of \(\Tetramod[B][][][]\)
is monoidally equivalent to the category \(\Tetramod[B][B][][B]\) of \emph{Hopf trimodules}\note{%
  The ``trimodule'' terminology is based on~\cite{shoikhet09:hopf}
  and was suggested by Ulrich Krähmer,
  in lieu of more unwieldy notations like ``\((2,1)\)\hyp{}Hopf module''.%
}:
\(B\)\hyp{}\(B\)\hyp{}bi\-co\-mo\-dules
with an additional left \(B\)\hyp{}action
that is a \(B\)\hyp{}\(B\)\hyp{}bicomodule morphism.

\begin{reptheorem}{thm:hopf-trimodules-are-lax-monoidal-functors}
  For a bialgebra \(B\) and \(\cat{V} \defeq \Tetramod[B]\),
  there is a monoidal equivalence
  \[
    \Trimod \simeq \mathsf{LexfLax}\cat{V}\mathsf{Mod}(\cat{V}, \cat{V})
  \]
  between the category of Hopf trimodules, and
  the category of left exact finitary lax \(\cat{V}\)\hyp{}module endofunctors on \(\cat{V}\).
\end{reptheorem}

Such structures feature prominently in the quasi-bialgebraic generalisation of the fundamental theorem of Hopf modules,
see~\cite{hausser99:integ-theor-quasi-hopf-algeb,saracco17:hopf}.
This equivalence matches a lax \(\Tetramod[B][][][]\)\hyp{}module monad with a Hopf trimodule algebra:
a Hopf trimodule \(A\) together with maps \(A \cotens A \to A\) and \(B \to A\) satisfying the usual associativity and unitality axioms for an algebra object.

If \(B\) is infinite-dimensional, \cref{mainnonsenseinjective}
should yield reconstruction in terms of an \emph{oplax} \(\Tetramod[B][][][]\)\hyp{}module comonad \(\cohom{X,\blank \triangleright X}\),
using the existence of injective objects in \(\Tetramod[B][][][]\).
Since our result gives an algebraic realisation only for the \emph{lax} \(\Tetramod[B][][][]\)\hyp{}module functors,
we restrict ourselves to the case where \(\blank \triangleright X\) admits both a left and a right adjoint,
by assuming \(\blank \triangleright X\) to be exact.
In that case, we obtain an adjunction \(\cohom{X,\blank \triangleright X} \dashv \hom{X,\blank \triangleright X}\).

As the monad is right adjoint,
the Eilenberg–Moore category of \(\cohom{X,\blank \triangleright X}\) is not equivalent to the category of modules over the associated Hopf trimodule algebra \(\hom{X,X}\),
but rather to the category of its \emph{contramodules}:
structures extensively studied in the setting of so-called semi-infinite homological algebra,~\cite{positselski10:homol}.
On the other hand, the Kleisli categories are equivalent,
and once again control the \(\Tetramod[B][][][]\)\hyp{}module structure,
which can thus be read off directly from the Hopf trimodule algebra.
We obtain the following result.

\begin{reptheorem}{finitetrihopfreconstruction}
  Let \(\cat{M}\) be a locally finite abelian \(\tetramodfd[B][][][]\)\hyp{}module category
  satisfying the assumptions of \cref{mainnonsenseinjective}.
  There exists a Hopf trimodule algebra \(A \in \tetramodfd[B][B][][B]\),
  such that there is an equivalence
  \(\Ind(\cat{M}) \simeq A\textnormal{-Contramod}\)
  of \(\tetramodfd[B]\)\hyp{}module categories,
  where the \(\tetramodfd[B]\)\hyp{}module structure on the right-hand side is extended from the category of free \(A\)\hyp{}module.
  This equivalence restricts to \(\cat{M} \simeq A\textnormal{-contramod}\).
\end{reptheorem}

Since this algebraic realisation of our reconstruction results may at first glance seem difficult to apply in calculations,
we give two explicit examples in which we determine a trimodule algebra for a given module category.

The first is \cref{coreps}, where \(B\) is the semigroup algebra for the unique two-element monoid which is not a group.
This example is intended to be very similar to~\cite[Example~2.20]{douglas19},
in that a simple object of \(\cat{C}\) acts as zero on an indecomposable object.
We also show that applying the ordinary,
``rigid'' reconstruction procedure on this example yields the same algebra object in \(\tetramodfd[B][][][]\)
as that corresponding to the regular action of \(\tetramodfd[B][][][]\) on itself,
and thus fails to classify module categories.

Our second example is given in \cref{fiberfunctorthings}.
Supposing \(B\) to be arbitrary,
we determine the Hopf trimodule algebra underlying the module functor
given by the fibre functor \(\tetramodfd[B] \to \tetramodfd\).

Finally, the Morita-theoretic results of \cref{injectivebijection} give us a precise notion of Morita equivalence for Hopf trimodule algebras
with respect to their contramodules in \cref{contramoritadef}.
This results in a Morita Theorem for \(\tetramodfd[B][][][]\)\hyp{}module categories, \cref{contramodulereuters}.

As another application, we use trimodule reconstruction
to give a categorical interpretation of a variant of the fundamental theorem of Hopf modules.

\begin{reptheorem}{LaxIsStrongThenLeftRigid,cor:twisted-anti-iff-bestoftimes}
  The functor \(\Tetramod[B] \to \Tetramod[B][B][][B]\) corresponds to
  the inclusion of strong \(\Tetramod[B]\)\hyp{}module endofunctors in the category of lax \(\Tetramod[B]\)\hyp{}module endofunctors,
  under the Yoneda lemma and \cref{thm:hopf-trimodules-are-lax-monoidal-functors}.
  The latter—and hence also the former—functor is an equivalence if and only if
  \(\tetramodfd[B]\) is left rigid,
  which is the case if and only if \(B\) has a twisted antipode.
\end{reptheorem}

In \cref{diffusion},
we show that the fusion operators of a Hopf monad \(T\) on a monoidal category \(\cat{W}\)
can be interpreted as coherence cells for the canonical oplax module structure on \(T\).
This gives a strong converse to the fact that (op)lax module functors over a rigid category are automatically strong,
by providing a distinguished module functor which is strong if and only if the category is rigid.
This result, interpreting additional structure of a monad on a monoidal category as morphisms for an (op)lax module functor structure on said monad,
is similar to the results of~\cite[Theorems~3.17 and~3.18]{flake24:froben-drinf},
which studies Frobenius monoidal functors rather than Hopf monads.

\section{Bicomodules and their graphical calculus}\label{sec:bicomodules}

\textsc{Before introducing Hopf trimodules},
let us first talk about bicomodules over \(B\)
in the sense of \cref{sec:algebra-and-module-objects}
and their string diagrammatic representation.

\begin{hypothesis}\label{hyp:hopf-trimodule}
  For this chapter, fix a field \(\Bbbk\) and a bialgebra \(B\) over it.
  Write \(\cat{V} \defeq \bbcomod\) for the category of left \(B\)\hyp{}comodules.
\end{hypothesis}

After introducing all of the relevant concepts and notation,
our first goal is to prove the following theorem.

\begin{theorem}\label{thm:hopf-trimodules-are-lax-monoidal-functors}
  There is a monoidal equivalence
  \begin{equation}\label{eq:hopf-trimodule-equiv}
    \begin{aligned}
      \Trimod &\to \mathsf{LexfLax}\cat{V}\mathsf{Mod}(\cat{V}, \cat{V})\\
      X &\mapsto (X \cotens \blank, \chi)
    \end{aligned}
  \end{equation}
  between the category of Hopf trimodules, and
  the category of left exact finitary lax \(\cat{V}\)\hyp{}module endofunctors on \(\cat{V}\),
  where \(\chi\) is the interchange morphism of \cref{def:bicomodule-interchange},
  see \cref{eq:hopf-trimod-to-lax-module-functors--functor-definition}
\end{theorem}

The quasi\hyp{}inverse of \cref{eq:hopf-trimodule-equiv}
evaluates a left exact finitary lax \(\cat{V}\)\hyp{}module functor at the injective generator \(B\).
We defer the proof of \cref{thm:hopf-trimodules-are-lax-monoidal-functors} until \cref{sec:hopf-trimodules-lax-mod-funs}.

\begin{notation}
  A bicomodule over \(B\) consists of a vector space \(X\),
  a left coaction \(\lambda \from X \to B \kotimes X\),
  and a right coaction \(\rho \from X \to X \kotimes B\), such that
  \[
    (B \kotimes \rho) \circ \lambda = (\lambda \kotimes B) \circ \rho.
  \]
  In graphical notation, we write
  \begin{equation}\label{eq:bicomodule-condition}
    \tikzfig{eq:bicomodule-condition}
  \end{equation}
\end{notation}

\begin{example}\label{ex:cotensor-product-of-bicomodules}
  Analogously to \cref{ex:cotensor-product-of-comodules},
  we can form the \emph{cotensor product} of two bicomodules \((X, \lambda^X, \rho^X)\) and \((Y, \lambda^Y, \rho^Y)\) over \(B\).
  This is the equaliser
  \[
    \begin{mytikzcd}[ampersand replacement=\&]
      {X \cotens Y} \& {X \kotimes Y} \& {X \kotimes B \kotimes Y.}
      \arrow[hook, from=1-1, to=1-2]
      \arrow["{\rho^X \kotimes Y}", shift left=1, from=1-2, to=1-3]
      \arrow["{{\raisebox{-0.8em}{\(X \kotimes \lambda^Y\)}}}"', shift right=1, from=1-2, to=1-3]
    \end{mytikzcd}
  \]
\end{example}

Our graphical calculus has to differentiate between
the tensor product of two bicomodules over \(\Bbbk\)
and their cotensor product.
In particular, we have to indicate which additional transformations are possible with the latter:
the equalised actions in the cotensor product will be annotated in grey.

\[
  \tikzfig{grey-cotensor}
\]\\
The fact that we may interchange the appropriate left and right actions will be indicated thusly:
\[
  \tikzfig{cotensor-interchange-actions}
\]

\section{From Hopf trimodules to lax module functors}\label{sec:hopf-trimodules-lax-mod-funs}

\textsc{In this chapter we give}
a proof of \cref{thm:hopf-trimodules-are-lax-monoidal-functors}.

\begin{definition}\label{def:hopf-tri-module}
  A \emph{Hopf trimodule} over \(B\) consists of a bicomodule \(X\)
  together with a left \(B\)-action \(\alpha\from B \kotimes X \to X\),
  such that \(\alpha\) is a left and right \(B\)\hyp{}comodule morphism.
  We write \(X \in \Trimod\).
\end{definition}

Alternatively, \cref{def:hopf-tri-module} could impose the conditions
that \(\lambda\) and \(\rho\) are left and right \(B\)\hyp{}module morphisms, respectively.
This equivalence is easily seen in a string diagrammatic reformulation:
\begin{equation}\label{eq:hopf-trimodule}
  \tikzfig{hopf-trimodule}
\end{equation}

\begin{definition}\label{def:bicomodule-interchange}
  Let \(X \in \Trimod\).
  For all \(M, N \in \bbcomod\), define the \emph{interchange morphism}
  \(\chi_{M, N} \from M \kotimes (X \cotens N) \to X \cotens (M \kotimes N)\)
  by\\\marginnote{%
      Notice the similarity of the interchange morphism
      with the Yetter--Drinfeld braiding of \cref{ex:yd-modules-are-braided}.%
    }
  \[
    \tikzfig{def:bicomodule-interchange}
  \]
\end{definition}

Observe in particular that
\[
  \tikzfig{rmk:interchange-unnessary-cotens}
\]
because the image of the coaction \(M \to B \kotimes M\) is contained in the cotensor product \(B \cotens M\).

\begin{lemma}\label{lem:interchange-properties}
  The interchange morphism
  is a well-defined left \(B\)\hyp{}comodule morphism.
\end{lemma}
\begin{proof}
  To show well-definedness,
  we need to show that the image of \(\chi_{M, N}\) is contained in \(X \cotens (M \kotimes N)\),
  for all \(M, N \in \bbcomod \).
  This follows by the calculation in \cref{fig:interchang-properties-calc1}.
  \begin{figure}[htbp]
    \centering
    \vspace{-0.3cm}\tikzfig{interchang-properties-calc1}\vspace{-0.3cm}% XXX
    \scaption[0.3cm]{%
      The image of \(\chi_{M, N}\) is contained in \(X \cotens (M \kotimes N)\).%
      \label{fig:interchang-properties-calc1}
    }
  \end{figure}

  Graphically, the fact that \(\chi\) is a left \(B\)\hyp{}comodule morphism means that
  \[
    \tikzfig{lem:interchange-properties-1}
  \]
  Using this, the calculation in \cref{fig:interchang-properties-calc2} finishes the proof.
  \begin{figure}[htbp]
    \centering
    \mathscale{0.85}{\tikzfig{interchang-properties-calc2}}\vspace{-0.3cm}% XXX
    \scaption[1.2cm]{%
      The arrow \(\chi\) is a morphism of left \(B\)\hyp{}comodules.%
      \label{fig:interchang-properties-calc2}%
    }
  \end{figure}
\end{proof}

\begin{lemma}\label{lem:interchange-is-braiding}
  The interchange morphism is a braiding; \ie,
  it is natural in both variables and the following diagrams commute
  for all \(M, N, P \in \bbcomod\):
  \begin{equation}\label[diagram]{eq:interchange-yb}
    \begin{mytikzcd}[ampersand replacement=\&,cramped,column sep=small]
      {M \kotimes (N \kotimes (X \cotens P))} \&\& {M \kotimes (X \cotens (N \kotimes P))} \\
      {(M \kotimes N) \kotimes (X \cotens P)} \& {X \cotens ((M \kotimes N) \kotimes P)} \& {X \cotens (M \kotimes (N \kotimes P))}
      \arrow["{M \kotimes \chi_{N, P}}", from=1-1, to=1-3]
      \arrow["{\chi_{M, N \kotimes P}}", from=1-3, to=2-3]
      \arrow["\alpha"', from=1-1, to=2-1]
      \arrow["\raisebox{-0.7em}{\(\chi_{M \kotimes N, P}\)}"', from=2-1, to=2-2]
      \arrow["\raisebox{-0.7em}{\(\alpha\)}"', from=2-2, to=2-3]
    \end{mytikzcd}
  \end{equation}
  \begin{equation}\label[diagram]{eq:interchange-unital}
    \begin{mytikzcd}[ampersand replacement=\&]
      {\Bbbk \kotimes (X \cotens M)} \&\& {X \cotens (\Bbbk \kotimes M)} \\
      \& {X \cotens N}
      \arrow["{\chi_{\Bbbk,N}}", from=1-1, to=1-3]
      \arrow["{X \cotens \lambda_N}", from=1-3, to=2-2]
      \arrow["{\lambda_{X \cotens M}}"', from=1-1, to=2-2]
    \end{mytikzcd}
  \end{equation}
\end{lemma}
\begin{proof}
  The interchange morphism is always natural in its second variable.
  To prove naturality in the first variable,
  let \(f \from M \to M'\) be a left comodule morphism.
  Then
  \[
    \tikzfig{lem:interchange-is-braiding}
  \]
  where the first equality follows by \(f\) being a left comodule morphism,
  and the second one is the naturality of the braiding.

  \Cref{eq:interchange-yb} commuting is equivalent to the following diagram,
  which is seen to be true by associativity of the action on \(X\).
  \[
    \tikzfig{lem:interchange-is-braiding-2}
  \]

  \Cref{eq:interchange-unital} follows by the unitality of the action on \(X\).
\end{proof}

\Cref{lem:interchange-properties,lem:interchange-is-braiding}
taken together say that the well-defined arrow
\[
  \chi\from \blank \kotimes (X \cotens \bblank) \nt X \cotens (\blank \kotimes \bblank)
\]
satisfies \cref{eq:interchange-yb,eq:interchange-unital}.
This, in turn, yields the following result.

\begin{proposition}\label{prop:interchange-comodule-functor}
  The pair \((X \cotens \blank, \chi)\) defines a lax \(\bbcomod\)\hyp{}module functor.
\end{proposition}

We can extend this correspondence to morphisms.

\begin{lemma}\label{lem:morphism-of-hopf-modules->module-transformation}
  Let \(f \in \Trimod(X, Y)\) be a morphism of Hopf trimodules.
  Then
  \[
    f \cotens \blank \from X \cotens \blank \nt Y \cotens \blank
  \]
  is a \(\bbcomod\)\hyp{}module transformation.
\end{lemma}
\begin{proof}
  We have to prove that
  \[
    \big(f \cotens (M \kotimes N)\big) \circ \chi^X_{M, N}
    =
    \chi^Y_{M, N} \circ \big(M \kotimes (f \cotens N)\big).
  \]
  In our graphical language, this means that
  \[
    \tikzfig{lem:morphism-of-hopf-modules->module-transformation}
  \]
  which follows immediately from \(f\) being a module morphism.
\end{proof}

Recall the notation \(\cat{V} \defeq \bbcomod \).
As a result of the previous considerations,
there exists a well-defined functor
\begin{equation}\label{eq:hopf-trimod-to-lax-module-functors--functor-definition}
  \Sigma \from \Trimod \to \mathsf{LexfLax}\cat{V}\mathsf{Mod}(\cat{V}, \cat{V}), \qquad X \mapsto (X \cotens \blank, \chi),
\end{equation}
To finish the proof of \cref{thm:hopf-trimodules-are-lax-monoidal-functors},
we have to show that \(\Sigma\) is monoidal, as well as an equivalence of categories.

\begin{lemma}\label{lem:cotens-of-hopf-modules-is-hopf-module}
  Let \(X,Y \in \Trimod\).
  Their cotensor product \(X \cotens Y\) is a Hopf trimodule,
  where the left action is given diagonally by \(b(x \otimes y) \defeq b_{(1)}x \otimes b_{(2)}y\).
\end{lemma}
\begin{proof}
  First, we show that the diagonal action of \(B\) is well-defined as a map \(B \kotimes (X \cotens Y) \to X \cotens Y\);
  \ie, that
  \[
    \tikzfig{lem:cotens-of-hopf-modules-is-hopf-module}
  \]
  This follows by the calculation in \cref{fig:cotens-of-hopf-modules-is-hopf-module-proof}.

  It is left to show that the action is a left and right \(B\)\hyp{}comodule morphism.
  The former case follows by \cref{fig:lem:cotens-of-hopf-modules-is-hopf-module-3},
  and in the latter case one calculates:
  \[
    \mathscale{0.90}{\tikzfig{lem:cotens-of-hopf-modules-is-hopf-module-4}}
  \]
  \begin{figure}[htbp]
    \centering
    \tikzfig{lem:cotens-of-hopf-modules-is-hopf-module-2}\vspace{-\onelineskip}% XXX
    \scaption[1cm]{%
      The diagonal action is well-defined on the cotensor product.%
      \label{fig:cotens-of-hopf-modules-is-hopf-module-proof}%
    }
  \end{figure}
  \begin{figure}[htbp]
    \centering
    \hrule\vspace{\onelineskip}%
    \tikzfig{lem:cotens-of-hopf-modules-is-hopf-module-3}\vspace{-\onelineskip}% XXX
    \scaption[-1cm]{%
      The action of the cotensor product is a right \(B\)\hyp{}comodule morphism.%
      \label{fig:lem:cotens-of-hopf-modules-is-hopf-module-3}%
    }
  \end{figure}
\end{proof}

\begin{proposition}\label{prop:braiding-composes-to-cotens}
  Let \(\chi^X\) and \(\chi^{Y}\) be as in \cref{def:bicomodule-interchange}.
  Then
  \[
    \chi^X \diamond \chi^Y \cong \chi^{X \cotens Y},
  \]
  where \(\diamond\) is the multiplicative cell for lax module morphisms;
  \ie,
  \[
    \big(X \cotens \blank, \chi^X\big) \diamond \big(Y \cotens \blank, \chi^Y\big) \defeq \big(X \cotens Y \cotens \blank, \chi^X \diamond \chi^Y\big).
  \]
  In other words, the functor \(\Sigma\)
  from \cref{eq:hopf-trimod-to-lax-module-functors--functor-definition}
  is strong monoidal.
\end{proposition}
\begin{proof}
  Associativity follows from coassociativity of the coaction and naturality of the braiding:
  \[
    \tikzfig{prop:braiding-composes-to-cotens}
  \]

  Unitality is immediate.
\end{proof}

To complete the proof of \cref{thm:hopf-trimodules-are-lax-monoidal-functors},
it is left to show that the strong monoidal functor \(\Sigma\)
from \cref{eq:hopf-trimod-to-lax-module-functors--functor-definition}
is an equivalence.
We will show that it is fully faithful and essentially surjective.

\begin{proposition}\label{prop:sigma-fully-faithful}
  The functor \(\Sigma\) is fully faithful.
\end{proposition}
\begin{proof}
  To show that \(\Sigma\) is faithful,
  suppose that \(f, g \from X \to Y\) are morphisms of Hopf trimodules,
  such that \(f \cotens \blank \cong g \cotens \blank\).
  Then
  \[
    \begin{mytikzcd}[ampersand replacement=\&]
      {X \cotens B} \& {Y \cotens B} \\
      X \& Y
      \arrow["{f \cotens B}", shift left=1, from=1-1, to=1-2]
      \arrow["g\cotens B"', shift right=1, from=1-1, to=1-2]
      \arrow["\cong"', from=1-1, to=2-1]
      \arrow["\cong", from=1-2, to=2-2]
      \arrow["f", shift left=1, from=2-1, to=2-2]
      \arrow["g"', shift right=1, from=2-1, to=2-2]
    \end{mytikzcd}
  \]
  commutes, so the result follows.

  To see that \(\Sigma\) is full,
  suppose that \(\varphi \from X \cotens \blank \nt Y \cotens \blank\) is a \(\cat{V}\)\hyp{}module transformation.
  By \cref{TakeuchiEilenbergWatts},
  we to show that the induced arrow
  \[
    \hat{\varphi}_B \from X \cong X \cotens B \xrightarrow{\ \varphi_B\ } Y \cotens B \cong Y
  \]
  is a morphism of Hopf trimodules.
  For a Hopf trimodule \(M \in \Trimod \),
  the morphism \(\varphi_M \from X \cotens M \to Y \cotens M\)
  and the induced morphism \(\hat{\varphi}_B\)
  can, in string diagrams,
  be depicted like this:
  \[
    \tikzfig{prop:sigma-fully-faithful}
  \]

  \cref{fig:sigma-phi-properties} collects several properties of \(\chi\) and \(\varphi\) in this notation.
  The arrow \(\hat{\varphi}_B\) is a left and right comodule morphism
  by the calculations in \cref{fig:sigma-phi-left-comodule-morphism}.
  Likewise, that \(\varphi_B\) is a left module morphism follows from:
  \[
    \tikzfig{prop:sigma-fully-faithful-2}
  \]
  \begin{figure}[htbp]
    \centering
    \tikzfig{sigma-phi-properties}
    \scaption{%
      Properties of  the arrows \(\chi\) and \(\varphi\).%
      \label{fig:sigma-phi-properties}%
    }
  \end{figure}
  \begin{figure}[htbp]
    \centering
    \tikzfig{sigma-phi-left-comodule-morphism}
    \scaption[-1cm]{%
      Verification that \(\hat{\varphi}_B\) is a left and right comodule morphism.%
      \label{fig:sigma-phi-left-comodule-morphism}
    }
  \end{figure}
\end{proof}

\newpage% XXX: Manual orphan prevention.
\begin{proposition}\label{prop:sigma-essentially-surjective}
  The functor \(\Sigma\) is essentially surjective.
\end{proposition}
\begin{proof}
  Suppose that \(\Phi \in \mathsf{LexfLax}\cat{V}\mathsf{Mod}(\cat{V},\cat{V})\).
  Since \(\Phi\) is left exact and finitary,
  by \cref{TakeuchiEilenbergWatts}
  there exists a \(B\)\hyp{}\(B\)\hyp{}bicomodule \(X\),
  such that \(\Phi \cong X \cotens \blank\) as functors.
  We will show that if \(X \cotens \blank\) is a lax module morphism,
  then \(X\) comes equipped with a left module structure,
  making it into a Hopf trimodule.

  Define a \(B\)\hyp{}action on \(X\) in the following way:
  \[
    \tikzfig{prop:sigma-essentially-surjective-1}
  \]
  Above, we have chosen to represent \(\chi_{B,B}\)%
  —see \cref{def:bicomodule-interchange}—%
  by a single box.
  The claim is that this action turns \(X\) into a Hopf trimodule.

  For brevity, we shall omit the grey action indicators in our string diagrammatic proofs,
  save when performing the respective transformation.

  To show that \(\alpha\) is a bicomodule morphism,
  it suffices to note that the left comodule morphism \(\chi_{B,B}\) is also a right comodule morphism,
  as \(\alpha\) is then composed of bicomodule morphisms.

  Notice that \(\chi\) is a natural transformation of functors which are left exact finitary in each variable, from
  \(\bbcomod \kotimes \bbcomod\) to \(\bbcomod\).
  Since \(\chi_{B,B}\) is the component of the transformation
  \[
    \chi_{B,\blank}\from B \kotimes (X \cotens \blank) \nt X \cotens (B \kotimes \blank)
  \]
  taken at the injective generator of \( \bbcomod \),
  it follows that it is a right comodule morphism;
  see \cref{prop:left exact-is-enough-on-injectives,TakeuchiEilenbergWatts}.
  Similarly, \(\chi_{B,B}\) is the component of \( \chi_{\blank,B}\from \blank \kotimes (X \cotens B) \nt X \cotens (\blank \kotimes B) \)
  taken at the injective generator,
  so \(\chi_{B,B}\) also intertwines this additional comodule structure:
  \[
    \tikzfig{prop:sigma-essentially-surjective-2}
  \]

  Let us now verify that \(\alpha\) is in fact a \(B\)\hyp{}module morphism.
  For the multiplicative identity of \(\alpha\), one calculates
  \[
    \tikzfig{prop:sigma-essentially-surjective-3}
  \]
  where (i) uses the fact that \((X \cotens \blank, \chi)\) is a \(\bbcomod\)\hyp{}module functor.

  The unitality of \(\alpha\) follows from the naturality of \(\chi\),
  the unital axiom of \((X \cotens \blank, \chi)\) being a \(\bbcomod\)\hyp{}module functor,
  and unitality of \(B\) and \(X\):
  \[
    \tikzfig{prop:sigma-essentially-surjective-4}
  \]

  It is left to show that \(\Sigma(\alpha) = \chi_{B,B}\).
  The necessary calculation is given in \cref{fig:sigma-essentially-surjective-hit},
  where equation (i) holds because \(\chi_{B,B}\) is a right comodule morphism in both variables,
  and (ii) follows by naturality of \(\chi\).
  \begin{figure}[hbtp]
    \centering
    \tikzfig{sigma-essentially-surjective-hit}\vspace{-\onelineskip}% XXX
    \scaption[1cm]{%
      Verification that \(\Sigma(\alpha) = \chi_{B,B}\).%
      \label{fig:sigma-essentially-surjective-hit}
    }
  \end{figure}
\end{proof}

\section{Contramodule reconstruction over Hopf trimodule algebras}\label{triconstruction}

\textsc{We stress that}
we are still working under \cref{hyp:deligne-reconstr,hyp:hopf-trimodule}.

\begin{definition}\label{def:hopf-trimodule-algebra}
  A \emph{Hopf trimodule algebra} consists of an algebra object in the monoidal category \(\big(\Trimod, \cotens, B\big)\).
  More explicitly, it consists of a Hopf trimodule \(A \in \Trimod\),
  together with trimodule morphisms \(\mu \from A \cotens A \to A\) and \(\eta \from B \to A\),
  satisfying associativity and unitality conditions.
\end{definition}

\begin{definition}\label{def:module-over-hopf-trimodule-algebra}
  A \emph{left module} over a Hopf trimodule algebra
  is a left \(A\)\hyp{}module object in the \(\Trimod\)\hyp{}module category \(\Tetramod[B]{}\)
  in the sense of \cref{ModuleInModule}.
  As such, it consists of a left \(B\)\hyp{}comodule \(M\)
  together with a \(B\)\hyp{}comodule homomorphism \(A \cotens M \to M\),
  satisfying associativity and unitality conditions with respect to \(\mu\) and \(\eta\).
\end{definition}

Similarly to \cref{def:module-over-hopf-trimodule-algebra} we define right modules over \(A\).
A module over \(A\) is \emph{free} if it is of the form \((A \cotens N, \mu \cotens N)\).

We can also dualise this definition:
we refer to objects of the category \(\Termod\) as \emph{Hopf termodules}.
A \emph{Hopf termodule coalgebra} is a coalgebra object \(C\) in the monoidal category \(\big(\Termod, \otimes_{B}, B\big)\),
and a \emph{comodule over \(C\)} is a left \(B\)\hyp{}module \(N\) together with a \(B\)\hyp{}module homomorphism \(N \to C \otimes_{B} N\),
satisfying usual coaction axioms.
Right comodules over \(C\) are defined analogously.

\begin{remark}
  Under the equivalence of \cref{eq:hopf-trimod-to-lax-module-functors--functor-definition}
  between \(\Trimod\) and \(\mathsf{LexfLax}\cat{V}\mathsf{Mod}(\cat{V}, \cat{V})\),
  the \(\Trimod\)\hyp{}module category \(\cat{V} = \Tetramod[B]{}\) corresponds to
  the evaluation action of \(\mathsf{LexfLax}\cat{V}\mathsf{Mod}(\cat{V}, \cat{V})\) on \(\cat{V}\).
  Thus, the category of modules over a Hopf trimodule algebra \(A\) is precisely
  the Eilenberg--Moore category for the monad \(A \cotens \blank\).
  The subcategory of free modules over \(A\) is
  the image of the Kleisli category for this monad
  under the canonical embedding
  \[
    {(\Tetramod[B]{})}_{A \cotens \blank} \xrightarrow{\ \ \iota\ \ } {(\Tetramod[B]{})}^{A \cotens \blank}.
  \]
\end{remark}

Let \(A \in \Trimod\) be a Hopf trimodule algebra
such that \({}^{B}A\) is quasi-finite,
meaning that the corresponding lax \(\Tetramod[B]{}\)\hyp{}module functor \(A \cotens \blank\) has a left adjoint \(\mathrm{cohom}_{B\blank}(A,\blank)\),
which is canonically a comonad by \cref{einerKleinerKleisliÄquivalenz}.

\begin{definition}\label{contramodules}
  A \emph{left contramodule over \(A\)} is an object in \({(\Tetramod[B]{})}^{\mathrm{cohom}_{B\blank}(A,\blank)}\).

  Explicitly, this means that a left contramodule is a left \(B\)\hyp{}comodule \(M\)
  together with a left \(B\)\hyp{}comodule homomorphism \(M \to \mathrm{cohom}_{B\blank}(A,M)\),
  endowing \(M\) with the structure of a comodule in \(\Tetramod[B]{}\) over \(\mathrm{cohom}_{B\blank}(A,\blank)\).
\end{definition}

A contramodule is \emph{free} if it lies in the image of the canonical embedding
\[
  {(\Tetramod[B]{})}_{\mathrm{cohom}_{B\blank}(A,\blank)} \xhookrightarrow{\ \ \iota\ \ } {(\Tetramod[B]{})}^{\mathrm{cohom}_{B\blank}(A,\blank)}.
\]
In other words, it is of the form \(\mathrm{cohom}_{B\blank}(A,M)\),
for a left \(B\)\hyp{}comodule \(M\),
and its coaction is of the form
\[
  \mathrm{cohom}_{B\blank}(A,M)
  \xrightarrow{\ \mu^{*}\ } \mathrm{cohom}_{B\blank}(A \cotens A, M)
  \cong \mathrm{cohom}_{B\blank}(A, \mathrm{cohom}_{B\blank}(A, M)).
\]
We denote the category of free contramodules by \(A\textnormal{-ContramodFree}\).

Similarly, the category \(C\textnormal{-Contramod}\) of left contramodules over a Hopf termodule coalgebra \(C\) is
the Eilenberg--Moore category \({(\Tetramod[][B]{})}^{\mathrm{Hom}_{B\blank}(C,\blank)}\)
for the monad \(\mathrm{Hom}_{B\blank}(C,\blank)\) right adjoint to the comonad \(C \otimes_{B} \blank\).
The free contramodules are defined analogously to the case of Hopf trimodules.

\begin{remark}
  Contramodules similar to those of \cref{contramodules} were studied extensively in more homological settings,
  see in particular~\cite{positselski10:homol}.
  Forgetting the left \(B\)\hyp{}module structure,
  a Hopf trimodule algebra yields a so-called \emph{semialgebra},
  and by forgetting the left \(B\)\hyp{}comodule structure,
  a Hopf termodule coalgebra yields a so-called \emph{bocs}—bimodule object, coalgebra structure.
\end{remark}

\begin{lemma}\label{freecontramodmodmodmod}
  Let \(A\) be a Hopf trimodule algebra.
  The action
  \[
    V \lact (A \cotens M) \defeq A \cotens (V \otimes M),
    \qquad
    \text{for all \(V \in \kVect\) and \(M \in \bcomod\)}
  \]
  endows the category \(A\textnormal{-ModFree}\) of free \(A\)\hyp{}modules
  with a canonical module category structure over \(\Tetramod[B]{}\).
  Similarly, \(A\textnormal{-ContramodFree}\) is a module category over \(\bbcomod\) by means of the action
  \(V \lact \mathrm{cohom}_{B\blank}(A,N) \defeq \mathrm{cohom}_{B\blank}(A,V \otimes N)\).
  These two \(\Tetramod[B]{}\)\hyp{}module categories are canonically equivalent.
\end{lemma}
\begin{proof}
  By definition, we have \(A\textnormal{-ModFree} = {(\Tetramod[B]{})}_{A \cotens \blank}\).
  Since \(A\) is a Hopf trimodule algebra,
  the monad \(A \cotens \blank\) is a lax \(\Tetramod[B]{}\)\hyp{}module monad,
  so the Kleisli category \({(\Tetramod[B]{})}_{A \cotens \blank}\) is a \(\Tetramod[B]{}\)\hyp{}module category,
  as described in \cref{laxKlCorrespondence},
  and it is easy to check that the assignments described in the lemma correspond precisely to those of \emph{ibid}.

  Similar considerations apply to \(A\textnormal{-ContramodFree}\),
  using the equality
  \[
    A\textnormal{-ContramodFree} = {(\Tetramod[B]{})}_{\mathrm{cohom}_{B\blank}(A,\blank)}
  \]
  and the fact that,
  since \(\mathrm{cohom}_{B\blank}(A,\blank)\) is the left adjoint of the lax \(\Tetramod[B]{}\)\hyp{}module monad \(A \cotens \blank\),
  it is an oplax \(\Tetramod[B]{}\)\hyp{}module comonad, by \cref{porism:doctrinal-module-adjunctions}.

  The asserted equivalence between these two \(\Tetramod[B]{}\)\hyp{}module categories
  follows directly from \cref{KleisliMonadCorrespondence}.
\end{proof}

Recall from~\cite[1.12]{takeuchi77:morit} that \(\mathrm{Cohom}_{B\blank}(A,\blank)\) is exact if and only if \(A\) is an injective left \(B\)\hyp{}comodule.

\begin{proposition}
  If \(\leftidx{^B\!}{A}{}\) is injective as a left \(B\)\hyp{}module,
  then the \(\bbcomod\)\hyp{}module category structure on \(A\textnormal{-ContramodFree}\) of \cref{freecontramodmodmodmod}
  extends to a \(\bbcomod\)\hyp{}module category structure on \(A\textnormal{-Contramod}\),
  in the sense of \cref{extendingalways}.
\end{proposition}

\begin{proof}
  For the comonadic dual of \cref{extendingalways},
  it suffices to notice that the comonad \(\mathrm{Cohom}_{B\blank}(A,\blank)\) is left exact,
  which is the case if and only if \(\leftidx{^{B}\!}{A}{}\) is injective, by~\cite[1.12]{takeuchi77:morit}.
\end{proof}

\begin{theorem}\label{infinitetrihopfreconstruction}
  Let \(B\) be a bialgebra
  and let \(\cat{M}\) be a locally finite abelian module category over \(\bcomod\),
  which admits a \(\bbcomod\)\hyp{}injective \(\bbcomod\)\hyp{}cogenerator \(X \in \Ind(\cat{M})\),
  such that \(\blank \lact X\from \bcomod \to \Ind(\cat{M})\) is exact and quasi-finite.
  Then there is a Hopf trimodule algebra \((A, \mu, \eta)\),
  such that \({}^{B}\!A\) is quasi-finite and injective,
  hence finitely cogenerated injective,
  such that there is an equivalence of \(\Tetramod[B]{}\)\hyp{}module categories
  \[
    \Ind(\cat{M}) \simeq A\textnormal{-Contramod},
  \]
  between \(\Ind(\cat{M})\) and the category of left \(A\)\hyp{}contramodules.
  The \(\bbcomod\)\hyp{}module structure on \(A\textnormal{-Contramod}\) is extended,
  in the sense of \cref{extendable},
  from the category \(A\textnormal{-ContramodFree}\) of free left \(A\)\hyp{}contramodules.

  This equivalence restricts to an equivalence \(\cat{M} \simeq A\textnormal{-contramod}\),
  between \(\cat{M}\) and the category of finite\hyp{}dimensional \(A\)\hyp{}contramodules.
\end{theorem}
\begin{proof}
  Let us again write \(\cat{V} \defeq \Tetramod[B]{}\).
  Since \(\blank \lact X\from \bcomod \to \Ind(\cat{M})\) is right exact,
  its extension to \(\Ind(\bcomod) \cong \cat{V}\) admits a right adjoint \(\hom{X,\blank}\) by \cref{saftrex}.
  Similarly, since \(\blank \lact X\from \bcomod \to \Ind(\cat{M})\) is left exact,
  so is its extension to \(\cat{V}\),
  which is also quasi-finite by assumption.
  Thus, by \cref{saftlex},
  the functor \(\blank \lact X\from \cat{V} \to \Ind(\cat{M})\) also has a left adjoint \(\cohom{X, \blank}\).
  We thus obtain an oplax \(\cat{V}\)\hyp{}module comonad \(\cohom{X, \blank \lact X}\) on \(\cat{V}\),
  and a right adjoint lax \(\cat{V}\)\hyp{}module monad \(\hom{X, \blank \lact X}\).

  By \cref{thm:hopf-trimodules-are-lax-monoidal-functors},
  we have an isomorphism \(\hom{X, \blank \lact X} \cong \hom{X, B \lact X} \cotens \blank\) of lax \(\cat{V}\)\hyp{}module monads,
  where \(\hom{X, B \lact X}\) is a Hopf trimodule algebra.
  Let us denote \(\hom{X, B \lact X}\) by \(A\).
  Since \(A \cotens \blank\) admits a left adjoint, \({}^{B}\!A\) is quasi-finite.

  By uniqueness of adjoints, an isomorphism \(\cohom{X, \blank \lact X} \cong \mathrm{cohom}_{B-}(A,\blank)\) follows.
  This is an isomorphism of \(\cat{V}\)\hyp{}module comonads,
  since the \(\cat{V}\)\hyp{}module comonad structure
  is lifted from the isomorphic oplax \(\cat{V}\)\hyp{}module monad structures on the respective right adjoints.

  By \cref{mainnonsenseinjective},
  we have an equivalence \(\Ind(\cat{M}) \simeq \cat{V}^{\cohom{X, \blank \,\lact\, X}}\) of \(\cat{V}\)\hyp{}module categories.
  In particular, the functor \(\cohom{X, \blank}\) is comonadic
  and hence exact, see \cref{thm:SmallBeck}.
  By assumption, so is its right adjoint \(\blank \lact X\),
  whence the comonad \(\cohom{X, \blank \lact X}\) is exact,
  which, by~\cite[1.12]{takeuchi77:morit} and since it is isomorphic to \(\mathrm{cohom}_{B-}(A,\blank)\),
  implies the injectivity of \({}^{B}\!A\).

  The \(\cat{V}\)\hyp{}module structure on \(\cat{V}^{\cohom{X,\blank \,\lact\, X}}\) is extended from \(\cat{V}_{\cohom{X,\blank \,\lact\, X}}\).
  By the isomorphism \(\cohom{X, \blank \lact X} \cong \mathrm{cohom}_{B-}(A,\blank)\) of the involved comonads,
  we have an equivalence of \(\cat{V}\)\hyp{}module categories
  \[
    \cat{V}^{\cohom{X, \blank \,\lact\, X}} \simeq \cat{V}^{\mathrm{cohom}(A,\blank)},
  \]
  and an equivalence \(\cat{V}_{\cohom{X, \blank \,\lact\, X}} \simeq \cat{V}_{\mathrm{cohom}(A,\blank)}\).
  Further, by definition
  \[
    \cat{V}^{\mathrm{cohom}(A,\blank)} = A\textnormal{-Contramod}
    \qquad\text{and}\qquad
    \cat{V}_{\mathrm{cohom}(A,\blank)} = A\textnormal{-ContramodFree},
  \]
  which yields the desired equivalence \(\Ind(\cat{M}) \simeq A\textnormal{-Contramod}\).

  The restricted equivalence \(\cat{M} \simeq A\textnormal{-contramod}\) follows from \cref{locomonads}.
  In particular, we use the observation that a contramodule is compact
  if and only if its underlying \(B\)\hyp{}comodule is such,
  if and only if it is finite\hyp{}dimensional.
\end{proof}

In the finite\hyp{}dimensional setting, a similar result uses projective objects.

\begin{theorem}\label{finitetrihopfreconstruction}
  Let \(B\) be a finite\hyp{}dimensional bialgebra
  and let \(\cat{M}\) be a finite abelian module category over \(\bmodule\),
  which admits an object \(X \in \cat{M}\) such that
  \(X\) is a \(\bmodule\)\hyp{}projective \(\bmodule\)\hyp{}generator; and
  \(\blank \lact X\from \bmodule \to \cat{M}\) is exact.

  Then there is a finite\hyp{}dimensional Hopf termodule coalgebra \(C \in \termod\),
  such that \({}_{B}C\) is projective and there is an equivalence of \(\bmodule\)\hyp{}module categories
  \[
    \cat{M} \simeq C\textnormal{-contramod}
  \]
  between \(\cat{M}\) and the category of finite\hyp{}dimensional left \(C\)\hyp{}contramodules.
  The \(\bmodule\)\hyp{}module structure on \(C\textnormal{-contramod}\)
  is extended from the free left \(A\)\hyp{}contramodules.
\end{theorem}
\begin{proof}
  This follows analogously to \cref{infinitetrihopfreconstruction},
  using \cref{mainnonsenseprojective} instead of \cref{mainnonsenseinjective}
  and \cref{finitemonads} instead of \cref{locomonads}.
\end{proof}

In the semisimple case, contramodules become equivalent to modules.

\begin{definition}
  We say that a Hopf trimodule algebra \(A\) is \emph{semisimple}
  if the monad \(A \cotens \blank\) is semisimple in the sense of \cref{def:semisimple-monad}.
\end{definition}

\begin{proposition}\label{semisimpletrimodulealgebra}
  For a semisimple Hopf trimodule algebra \(A \in \Trimod\),
  there is an equivalence \(\lMod{A} \simeq A\textnormal{-Contramod}\) of module categories over \(\Tetramod[B]{}\).
\end{proposition}
\begin{proof}
  This is \cref{ModuleEMSemisimple} applied to the monad \(A \cotens \blank\).
\end{proof}

\subsection{Morita equivalence for contramodules}

\begin{definition}~\label{contramoritadef}
  We say that two Hopf trimodule algebras \(A, A' \in \trimod\),
  such that \({}^{B}\!A\) and \({}^{B}\!A'\) are quasi-finite,
  are \emph{contraMorita equivalent} if the comonads \(\mathrm{cohom}_{B\blank}(A,\blank)\) and \(\mathrm{cohom}_{B\blank}(A',\blank)\)
  are Morita equivalent.
\end{definition}

More explicitly, \cref{contramoritadef} says that \(A\) and \(A'\) are contraMorita equivalent
if there are \(M, N \in \Tetramod[B][B][][B]\) with coaction morphisms in \(\Tetramod[B][B][][B]\)
\[
  \begin{aligned}
    &\mathrm{la}_{M,A}\colon M \to \mathrm{cohom}_{B\blank}(A,M),
    &&\mathrm{ra}_{N,A}\colon N \to N\cotens_{B} \mathrm{cohom}_{B\blank}(A,B), \\
    &\mathrm{la}_{N,A'}\colon N \to \mathrm{cohom}_{B\blank}(A',N),
    &&\mathrm{ra}_{M,A'}\colon M \to M\cotens_{B} \mathrm{cohom}_{B\blank}(A',B), \\
  \end{aligned}
\]
such that
\begin{itemize}
  \item
  there is an isomorphism in \(\Tetramod[B][B][][B]\) between the equaliser of
  \[
    \mathscale{0.91}{\begin{mytikzcd}[ampersand replacement=\&]
	\& {M \cotens_{B}\mathrm{cohom}_{B\blank}(A',B)\cotens_{B} N} \\
	{M \cotens_{B} N} \&\& {M \cotens_{B} \mathrm{cohom}_{B\blank}(A',N)}
	\arrow["{M \cotens_{B} \partial_{A',B,N}}"{pos=0.7}, from=1-2, to=2-3]
	\arrow["\simeq"', from=1-2, to=2-3]
	\arrow["{\mathrm{ra}_{M,A'} \cotens_{B} N}"{pos=0.3}, from=2-1, to=1-2]
	\arrow[from=2-1, to=2-3]
      \end{mytikzcd}}
  \]
  and \(A\),
  intertwining the two coactions \(A \to A \cotens_{B} \mathrm{cohom}_{B-}(A,B)\) and \(A \to \mathrm{cohom}_{B-}(A,A)\)
  coming from the algebra structure on \(A\),
  with the coactions on the equaliser induced by \(\mathrm{ra}_{N,A}\) and \(\mathrm{la}_{M,A}\).
  Here, \(\partial_{A',B,N}\) is the isomorphism of~\cite[1.13]{takeuchi77:morit},
  which is invertible because \({}^{B}\!A'\) is injective.
  \item a similar isomorphism in \(\Tetramod[B][B][][B]\)
  between a similarly defined equaliser reversing the roles of \(M\) and \(N\), and \(A'\) exists.
\end{itemize}

\begin{proposition}\label{contraspberry}
  For \(A,A'\) as in \cref{contramoritadef}, the following are equivalent:
  \begin{thmlist}[noitemsep]
    \item\label{contrabanda} \(A\) and \(A'\) are contraMorita equivalent,
    \item\label{bigequivatti} \(A\textnormal{-Contramod} \simeq A'\textnormal{-Contramod}\) as \(\Tetramod[B]{}\)\hyp{}module categories, and
    \item\label{piccoloequivatti} \(A\textnormal{-contramod} \simeq A'\textnormal{-contramod}\) as \(\tetramodfd[B]{}\)\hyp{}module categories.
  \end{thmlist}
\end{proposition}
\begin{proof}
  The equivalence of~\ref{contrabanda} and~\ref{bigequivatti} is an instance of \cref{Moritabybimodules},
  for \(\cat{C} = \Tetramod[B]{}\).
  Since an equivalence of categories preserves compact objects,~\ref{piccoloequivatti} follows from~\ref{bigequivatti}.
  Finally, if \(F\) is an equivalence as in~\ref{piccoloequivatti}, then \(\Ind(F)\) is an equivalence as in~\ref{bigequivatti}.
\end{proof}

Combining \cref{contraspberry} with \cref{infinitetrihopfreconstruction},
we find the following algebraic realisation of \cref{injectivebijection}:

\begin{theorem}\label{contramodulereuters}
  There is a bijection
  \[
    \faktor{\{\,(\cat{M},X) \text{ as in Theorem~\ref{infinitetrihopfreconstruction}}\,\}}{\cat{M} \simeq \cat{N}}
    \xleftrightarrow{\cong}
    \left\{
    \begin{aligned}
      &\text{Left finitely cogenerated} \\
      &\text{injective Hopf trimodule} \\
      &\text{algebras over \(B\)}
    \end{aligned}
    \right\}
    \!\!\Big/\!\simeq_{\text{cMorita}}\!.
  \]
\end{theorem}

\section{A semisimple example of non-rigid reconstruction}\label{coreps}

\textsc{We now give a concrete application} of \cref{mainnonsenseinjective,infinitetrihopfreconstruction},
describing module categories for non-rigid categories
where the ordinary reconstruction procedure for rigid categories would not yield the correct result.

Let \(S = \{e,s\}\) be the commutative two-element monoid which is not a group,
with identity element \(e\).
In particular, \(s^{2} = s\) is an idempotent.
The category \(\mathsf{rep}(S)\) of finite\hyp{}dimensional modules over the monoid algebra \(\Bbbk[S]\) is semisimple,
with two isoclasses of simple objects given by \(1\)\hyp{}dimensional representations \(\Bbbk_{\mathrm{triv}}\) and \(\Bbbk_{\mathrm{grp}}\).
In \(\Bbbk_{\mathrm{triv}}\), the element \(s\) acts by \(1\),
while in \(\Bbbk_{\mathrm{grp}}\) it acts by \(0\).
The standard comultiplication on \(\Bbbk[S]\), in which \(e\) and \(s\) are group-like, makes turns \(\Bbbk[S]\) into a bialgebra.
The category \(\mathsf{corep}(S)\) of finite\hyp{}dimensional comodules over \(\Bbbk[S]\) is equivalent to that of \(S\)\hyp{}graded finite\hyp{}dimensional spaces, \(\kvect^{S}\).
Thus, it too is semisimple of rank two—we denote the simple objects by \(\delta_{e}\) and \(\delta_{s}\).

The monoidal structure on \(\mathsf{rep}(S)\) satisfies \(\Bbbk_{\mathrm{grp}} \otimes \Bbbk_{\mathrm{grp}} \cong \Bbbk_{\mathrm{grp}}\).
In fact, this presentation of the Grothendieck ring \([\mathsf{rep}(S)]\) of \(\mathsf{rep}(S)\) determines \(\mathsf{rep}(S)\) uniquely as a semisimple monoidal category, see~\cite[Proposition~3.5]{sun23:notes}.
Thus, we find monoidal equivalences
\[
  \mathsf{rep}(S) \simeq \mathsf{corep}(S) \simeq \kvect^{S},
\]
since in the latter case we have \(\delta_{s} \otimes \delta_{s} \cong \delta_{s}\).

Let us now focus on the coalgebraic setup of \cref{mainnonsenseinjective},
studying the module categories over \(\mathsf{corep}(S)\).

A semisimple \(\mathsf{corep}(S)\)\hyp{}module category \((\cat{M}, \lact)\) is indecomposable\note{%
  A \(\cat{C}\)\hyp{}module category \(\cat{M}\) is called \emph{indecomposable} if
  it does not split as a direct sum of \(\cat{C}\)\hyp{}module subcategories.%
}
if and only if
the matrix \({[\delta_{s}]}_{\cat{M}}\)
describing the action of \(\delta_{s}\) on the Grothendieck ring \([\cat{M}]\) in the basis of simple objects for \(\cat{M}\)
is indecomposable.
But \([\delta_{s}]\) is idempotent, so it must be the \(1\times 1\)\hyp{}matrix \((1)\), and,
as a category, we must have \(\cat{M} \simeq \kvect\).
We then either have \(\delta_{s} \lact \blank \cong \Id_{\cat{M}}\),
or \(\delta_{s} \lact \blank \cong 0\).
One can show that any two module categories satisfying one of these conditions are equivalent,
by showing that the second cohomology group \(H^{2}(S, \Bbbk^{\times})\) is trivial,
similarly to~\cite[Example~7.4.10]{Etingof2015}.
Indeed, \(\psi \in Z^{2}(S, \Bbbk^{\times})\) must satisfy the relation \(\psi(e,s) = \psi(e,e) = \psi(s,e)\),
and thus can be written as \(\psi(x,y) = \sigma(x) - \sigma(xy) + \sigma(y)\),
by setting \(\sigma(e) \defeq \psi(e,e)\) and \(\sigma(s) \defeq \psi(s,s)\).
The condition \(\delta_{s} \lact \blank \cong \Id_{\cat{M}}\) is satisfied by the fibre functor
\(\mathrm{Fr}\from \mathsf{corep}(S) \to \kvect\).

Assume instead that we have an indecomposable module category \(\cat{M}\) satisfying \(\delta_{s} \lact \blank \cong 0\).
Let \(X \in \cat{M}\) be a simple object.
One can verify that \(X\) satisfies the assumptions of \cref{mainnonsenseinjective}.
Given \(V \in \mathsf{corep}(S)\), we have
\[
  \lceil X,V \lact X \rceil_{s} = {\mathsf{corep}(S)}(\lceil X, V\lact X \rceil, \delta_{s}) \cong {\mathsf{corep}(S)}(V \lact X, \delta_{s} \lact X),
\]
and similarly for \(\delta_{e}\).

Using \(\delta_{s} \lact X \cong 0\) and \(\delta_{e} \lact X \cong X\),
we find that \(\lceil X, \delta_{e} \lact X \rceil \cong \delta_{e}\) and \(\lceil X, \delta_{s} \lact X \rceil = 0\).
To describe the bicomodule corresponding to the functor \(\lceil X, \blank \lact X \rceil\)
under the correspondence of \cref{TakeuchiEilenbergWatts},
we observe that a bicomodule can be identified with an \(S \times S\)\hyp{}graded space,
and the cotensor product is given by
\[
  {(V \cotens W)}_{x,y} = \bigoplus_{z \in S} V_{x,z} \otimes W_{z,y}.
\]
Using this, we find that \(\lceil X, \blank \lact X \rceil \cong \delta_{e,e} \cotens \blank\),
where \(\delta_{e,e}\) is the one-dimensional space concentrated in degree \((e,e)\).
The unit Hopf trimodule homomorphism \(\Bbbk[S] \to \lceil X, \Bbbk[S] \lact X \rceil\)
must thus have the submodule \(\Bbbk\{s\}\) as its kernel,
so the reconstructing Hopf trimodule \(\lceil X, \Bbbk[S] \lact X \rceil\) is given by \(\Bbbk[S]/\Bbbk\{s\}\).
As a bicomodule, it is indeed concentrated in degree \((e,e)\),
and as a module it is isomorphic to \(\Bbbk_{\mathrm{grp}}\).
The composition
\[
  \lceil X, \delta_{e} \lact X\rceil \otimes \lceil X, \delta_{e} \lact X\rceil \to \lceil X, \delta_{e} \lact X\rceil
\]
sending \(\mathrm{id}_{X} \otimes \mathrm{id}_{X}\) to \(\mathrm{id}_{X}\) corresponds to the algebra structure
\[
  \begin{aligned}
    \Bbbk[S]/\Bbbk\{s\} \cotens \Bbbk[S]/\Bbbk\{s\} \simeq \Bbbk[S]/\Bbbk\{s\} \otimes \Bbbk[S]/\Bbbk\{s\} &\to \Bbbk[S]/\Bbbk\{s\} \\
    (e + \Bbbk\{s\}) \otimes (e + \Bbbk\{s\}) &\mapsto (e + \Bbbk\{s\}),
  \end{aligned}
\]
and this multiplication defines a Hopf trimodule algebra.

To check that the category \(\Bbbk[S]/\Bbbk\{s\}\textnormal{-Contramod}\)
satisfies the properties we have assumed of \(\cat{M}\) previously,
note that the free module \(\Bbbk[S]/\Bbbk\{s\} \cotens B\) is one-dimensional,
hence semisimple,
showing that the category of modules over \(\Bbbk[S]/\Bbbk\{s\}\) is semisimple of rank one.
Thus, by \cref{semisimpletrimodulealgebra} we have
\[
  \Bbbk[S]/\Bbbk\{s\}\textnormal{-Contramod} \simeq \Bbbk[S]/\Bbbk\{s\}\textnormal{-Mod}.
\]
Further, \(\delta_{s} \lact \Bbbk[S]/\Bbbk\{s\} = \Bbbk[S]/\Bbbk\{s\} \cotens \delta_{s} = 0\), which corresponds to \(\delta_{s} \lact X = 0\).

If we instead follow the reconstruction procedure described in~\cite[Chapter~7]{Etingof2015},
the reconstructing algebra object would be the object \(A\) of \(\mathsf{corep}(S)\)
representing the presheaf \(\cat{M}(\blank \lact X, X)\).
We see that we must have \(A \cong \delta_{e}\),
and the algebra structure on \(\delta_{e}\) is unique.
However, \(\delta_{e} \cong 1_{\mathsf{corep}(S)}\) and so
\[
  \mathrm{Mod}(\delta_{e}) \simeq \mathsf{corep}(S) \not\simeq \cat{M},
\]
showing that the reconstruction procedure for module categories over rigid monoidal categories fails in this non-rigid case.
It also illustrates that this failure can be observed already in the semisimple case.

\section{A Hopf trimodule algebra constructing the fibre functor}\label{fiberfunctorthings}

\textsc{This section discusses the general construction}
of a Hopf trimodule algebra,
which covers the remaining indecomposable semisimple module category over the category \(\mathsf{corep}(S)\) considered in \cref{coreps}.

\begin{proposition}\label{prop:B-dot-B-trimodule}
  The \(\Bbbk\)\hyp{}vector space \(B \kotimes B\) forms a Hopf trimodule
  with coactions
  \[
    b \otimes c \mapsto b_{(1)} \otimes b_{(2)} \otimes c,\qquad\qquad
    c \otimes b \mapsto c \otimes b_{(1)} \otimes b_{(2)},
  \]
  and action
  \[
    x \otimes b \otimes c \mapsto x_{(1)}b \otimes x_{(2)}b.
  \]
  We shall denote it by \(B \bullet B\).
\end{proposition}
\begin{proof}
  It is clear that \(B \bullet B\) is a bicomodule,
  as both of the coactions are given by that of \(B\) on the respective Let.

  Thus, we first verify that the action is a morphism of left and right comodules;
  \ie, that \(B \bullet B\) is in \(\Tetramod[][B][][B]\) and \(\Tetramod[B][B]\).
  The former follows by \cref{fig:b-bullet-b-left-mod-right-comod},
  and the latter is due to the calculation in \cref{fig:b-bullet-b-left-mod-left-comod}.
  \begin{figure}[htbp]
    \centering
    \tikzfig{b-bullet-b-left-mod-right-comod}\vspace{-\onelineskip}
    \scaption[1cm]{%
      The action of \(B \bullet B\) is a morphism of left comodules.%
      \label{fig:b-bullet-b-left-mod-right-comod}%
    }
  \end{figure}
  \begin{figure}[htbp]
    \centering
    \hrule\vspace{\onelineskip}
    \tikzfig{b-bullet-b-left-mod-left-comod}\vspace{-\onelineskip}
    \scaption[-.5cm]{%
      The action of \(B \bullet B\) is a morphism of right comodules.%
      \label{fig:b-bullet-b-left-mod-left-comod}
    }
  \end{figure}
\end{proof}

\begin{proposition}\label{prop:b-dot-b-is-algebra-object}
  The Hopf trimodule \(B \bullet B\) from \cref{prop:B-dot-B-trimodule} is an algebra object in \(\Tetramod[B][B][][B]\).
\end{proposition}
\begin{proof}
  We define the multiplication and unit of \(B \bullet B\) by
  \[
    \tikzfig{b-bullet-b-algebra-structure-def}
  \]
  where we have denote the counit \(\varepsilon\) of \(B\) with a small white dot,
  and have employed the same numbering system for representing the cotensor product as in \cref{sec:bicomodules}.
  It is clear that \(\mu\) is associative,
  left unitality follows from
  \[
    \tikzfig{b-bullet-b-mu-unitality}
  \]
  and right unitality is similar.

  One immediately has that \(\mu\) is a \(\Tetramod[B][][][B]\)\hyp{}morphism,
  and for \(\eta = \Delta\) this follows from coassociativity.
  To finish the proof we need only show that both are morphisms of left \(B\)\hyp{}modules,
  which follows from \cref{fig:b-bullet-b-module-morphisms}
  and the following calculation:
  \[
    \tikzfig{b-bullet-b-module-morphisms-2}
  \]
  \begin{figure}[htbp]
    \centering
    \tikzfig{b-bullet-b-module-morphisms}\vspace{-\onelineskip}% XXX
    \scaption{%
      The unit \(\eta = \Delta\)is a morphism of left \(B\)\hyp{}modules.%
      \label{fig:b-bullet-b-module-morphisms}%
    }
  \end{figure}
\end{proof}

\begin{definition}\label{jequivalence}
  Let \(J\from (B \bullet B)\textnormal{-ModFree} \to \kVect\) be the functor given by \(J((B \bullet B) \cotens M) = M\) and by local maps
  \[
    \begin{aligned}
      J_{M,N}&\from (B\bullet B)\textnormal{-ModFree}\big((B\bullet B)\cotens M,  (B\bullet B)\cotens N \big) \to \mathrm{Hom}_{\Bbbk}(M,N) \\
      \sigma &\mapsto (\varepsilon \otimes \varepsilon \cotens \mathrm{id}_{N})\circ(\sigma(1 \otimes 1 \otimes \blank))
    \end{aligned}
  \]
\end{definition}

\begin{lemma}
  \Cref{jequivalence} yields an equivalence of \(\bbcomod\)\hyp{}module categories.
\end{lemma}
\begin{proof}
  First, assume that \(J\) is a functor;
  we will show it is an equivalence of (\(\bbcomod\))\hyp{}module categories.
  Let \(V \in \kVect\).
  Essential surjectivity follows by
  \[
    V \cong \Bbbk^{\dim{V}} \cong J((B\bullet B)\cotens \Bbbk_{\mathrm{triv}}^{\dim{V}}).
  \]

  The fact that \(J\) is fully faithful follows by
  \(J_{M,N}\) being defined as the composite of the isomorphisms
  \begin{align*}
    (B \bullet B)&\textnormal{-ModFree}\big((B\bullet B)\cotens M,  (B\bullet B)\cotens N\big)
             \xrightarrow{\blank \circ \eta} \bbcomod(M,  (B\bullet B)\cotens N) \\
           & \xrightarrow{(B \otimes \varepsilon \cotens N) \circ \blank} \bbcomod(M,  B\otimes N)
             \xrightarrow{(\varepsilon \otimes N) \circ \blank} \kvect(M,N).
  \end{align*}
  The first morphism is an isomorphism of the form \(\cat{A}^T(TX, TY) \iso \cat{A}(X,TY)\),
  for the monad \(T \defeq (B \bullet B) \cotens \blank\) on the category \(\cat{A} \defeq \bbcomod\),
  see \cref{rmk:em-iso-and-so-on}.
  The second isomorphism is simply postcomposition with the isomorphism \((B \bullet B) \cotens N = B \otimes  B \cotens N \xrightarrow{B \otimes \varepsilon \cotens N} B \otimes N\),
  since \(B \cotens N \xrightarrow{\varepsilon \cotens N} N\) is an isomorphism.
  The third is an isomorphism \(\cat{B}^C(CX,CY) \iso \cat{B}(CX, Y)\),
  for the comonad \(K \defeq B \otimes \blank\) on the category \(\cat{B} \defeq \kVect\).

  Lastly, \(J\) is a (\(\bbcomod\))\hyp{}module functor:
  for \(W \in \bbcomod\), we have
  \[
    W \lact J(B\bullet B \cotens M) = W \otimes M = J(B\bullet B \cotens (W \otimes M)) = J(W \lact B \bullet B \cotens M).
  \]
  Thus, if \(J\) defines a functor, it is an equivalence of (\(\bbcomod\))\hyp{}module categories.

  It is left to verify the functoriality of \(J\).
  Let
  \begin{align*}
    \sigma &\in (B \bullet B)\text{-ModFree}((B \bullet B) \cotens M, (B \bullet B) \cotens N),\\
    \tau &\in (B \bullet B)\text{-ModFree}((B \bullet B) \cotens N, (B \bullet B) \cotens W).
  \end{align*}
  We represent these morphisms graphically by
  \[
    \tikzfig{represent-the-morphism-graphically-by}
  \]\\[-0.9em]

  Functoriality of \(J\) follows by \cref{fig:jequivalence-functoriality};
  (i) and (ii) hold due to
  \[
    \tikzfig{i-and-ii-hold-due-to}
  \]
  Equality (iii) is satisfied by \(\tau\) being a \((B \bullet B)\)\hyp{}module morphism,
  (iv) follows by the following calculation, using that \(B\) is a bialgebra:
  \[
    \varepsilon(b_{(1)}b_{(2)}) \otimes b_{(3)}
    = \varepsilon(b_{(1)})\varepsilon(b_{(2)}) \otimes b_{(3)}
    = \varepsilon(b_{(1)}\varepsilon(b_{(2)})) \otimes b_{(3)}
    = \varepsilon(b_{(1)}) \otimes b_{(2)}
    = b,
  \]
  and (v) holds by an analogous computation.
  \begin{figure}[htbp]
    \centering
    \tikzfig{jequivalence-functoriality}\vspace{-\onelineskip}% XXX
    \scaption[.5cm]{%
      The assignment \(J\) is functorial.%
      \label{fig:jequivalence-functoriality}%
    }
  \end{figure}
\end{proof}

\begin{corollary}
  The monad \((B \bullet B) \cotens \blank\) is semisimple, and there is an equivalence of \(\bbcomod\)\hyp{}module categories \((B \bullet B)\textnormal{-Mod} \simeq \kVect\).
\end{corollary}

\section{The fundamental theorem of Hopf modules}\label{sec:Hausser-Nill}

\begin{definition}[{\cite[Proposition~1.5.10]{montgomery93:hopf}}]
  Let \(B\) be a bialgebra.
  A \emph{twisted antipode} for \(B\) is an antipode for the bialgebra \(B^{\mathrm{cop}}\).\note{%
    Equivalently, a twisted antipode is an antipode for \(B^{\op}\).%
  }
\end{definition}

\begin{remark}
  A Hopf algebra admits a twisted antipode \(\tilde{S}\) if and only if its own antipode \(S\) is invertible,
  in which case we have \(S^{-1} = \tilde{S}\).
  Notably, a finite\hyp{}dimensional Hopf algebra always admits a twisted antipode.
\end{remark}

\Cref{antipodeleftrigid} below is a well-known result, see for example~\cite[Theorem~5.4.1.(2)]{Etingof2015}.
We emphasise that while its former part is seemingly opposite to the claim of~\cite[Theorem~5.4.1.(2)]{Etingof2015},
identifying the presence of an antipode with right, rather than left, rigidity, this is because we consider left, rather than right, comodules.
The latter claim of \cref{antipodeleftrigid} follows from the former by using the equivalence \({(\Tetramod[B])}^{\tensorop} \simeq \Tetramod[B^{\op}]\).

\begin{lemma}\label{antipodeleftrigid}
  Let \(B\) be a bialgebra.
  Then \(\tetramodfd[B]\) is right rigid if and only if \(B\) admits an antipode.
  Similarly, \(\tetramodfd[B]\) is left rigid if and only if \(B\) admits a twisted antipode.
\end{lemma}

The next theorem was shown
in the setting of quasi-bialgebras
in~\cite[Proposition~3.11]{hausser99:integ-theor-quasi-hopf-algeb}
and~\cite[Theorem~3.9]{saracco17:hopf},
as a generalisation of the structure theorem for Hopf modules of Larson and Sweedler~\cite{larson69:hopf}.

\begin{theorem}\label{HausserNill}
  Let \(B\) be a bialgebra.
  The functor \(\blank \otimes B\from \Tetramod[][B] \to \Tetramod[][B][B][B]\) is an equivalence if and only if \(B\) admits an antipode.
\end{theorem}

In the rest of this section we show that a variant of \cref{HausserNill} can be deduced
from the following general statement.

\begin{proposition}\label{LaxIsStrongThenLeftRigid}
  Let \(\cat{C}\) be a left closed monoidal category such that every left \(\cat{C}\)\hyp{}module endofunctor of \(\cat{C}\) is strong%
  —in other words,that the monoidal embedding
  \[
    \mathsf{Str}\cat{C}\mathsf{Mod}(\cat{C},\cat{C}) \longhookrightarrow \mathsf{Lax}\cat{C}\mathsf{Mod}(\cat{C},\cat{C})
  \]
  is an equivalence.
  Then \(\cat{C}\) is left rigid.
\end{proposition}
\begin{proof}
  Since \(\cat{C}\) is left closed,
  for every \(x \in \cat{C}\) there is a right adjoint \(\hom{x, \blank}\) to the strong \(\cat{C}\)\hyp{}module endofunctor \(\blank \otimes x\).
  By \cref{porism:doctrinal-module-adjunctions},
  the functor \(\hom{x, \blank}\) is a lax \(\cat{C}\)\hyp{}module endofunctor of \(\cat{C}\),
  and thus, by assumption, it is a strong module endofunctor.
  \Cref{prop:bicategorical-yoneda-correspondence-monads-and-algebras} yields an object \(\ld{\!{}x}\)
  such that there is an isomorphism of \(\cat{C}\)\hyp{}module functors \(\hom{x, \blank} \cong \blank \otimes \ld{\!{}x}\).
  Thus we obtain an adjunction of \(\cat{C}\)\hyp{}module functors \(\blank \otimes x \dashv \blank \otimes \ld{\!{}x}\),
  which by~\cite[102-103]{kelly72:many} shows that \((\ld{\!{}x}, x)\) is a dual pair in \(\cat{C}\).
  It follows that \(\cat{C}\) is left rigid.
\end{proof}

\begin{remark}\label{rmk:rigid-iff-lax-is-strong}
  Given a rigid monoidal category \(\cat{C}\),
  notice that every lax \(\cat{C}\)\hyp{}module endofunctor on \(\cat{C}\) is automatically strong by \cref{prop:lax-is-strong-when-rigid}.
  Together with \cref{LaxIsStrongThenLeftRigid},
  this yields a characterisation when a (left) closed monoidal category is (left) rigid,
  solely based on its lax \(\cat{C}\)\hyp{}module functors.
  This is a special property of module functors,
  as in \cref{sec:negative} we have already seen that ordinary adjunctions do not suffice for such a characterisation.
\end{remark}

We now define the functor \(B \otimes \blank\) analogous to \(\blank \otimes B\) of \cref{HausserNill}.

\begin{definition}\label{bestoftimes}
  Define the functor \(B \otimes \blank\from \Tetramod[B][][][] \to \Tetramod[B][B][][B]\)
  by sending a left \(B\)\hyp{}comodule \(M\) to the trimodule \(B \otimes M\) whose left and right coaction,
  as well as left action, is given by
  \[
    b \otimes m \mapsto b_{(1)}m_{(-1)} \otimes b_{(2)} \otimes m_{(0)}; \quad
    b \otimes m \mapsto b_{(1)}\otimes m \otimes b_{(2)}; \quad
    a \otimes b \otimes m \mapsto ab \otimes m;
  \]
  where we extends to morphisms in the natural way.
  It is easy to verify that \(B \otimes M\) defines a Hopf trimodule and also that \(B \otimes \blank\) defines a functor.
\end{definition}

\begin{remark}
  Comparing \cref{bestoftimes} with \cref{HausserNill}, we remark on two differences.
  First, the domain of the functor we study is the category of comodules over \(B\), rather than modules.
  This is merely to maintain consistence with \cref{sec:hopf-trimodules-lax-mod-funs},
  and a variant for modules can be formulated without greater difficulty.
  Second, the functor we study endows a left comodule with an additional comodule structure on the right,
  and a module structure on the left, rather than on the right.
  We will see that it is related to the functor \(\Tetramod{} \to \Tetramod[][B][][B]\),
  which is known to be an equivalence if \(B\) admits a twisted antipode;
  see for example~\cite[Section~1.9.4]{montgomery93:hopf}.
\end{remark}

We now give a categorical interpretation of the functor of \cref{bestoftimes}.
Recall the notation \(\cat{V} \defeq \Tetramod[B]\) of \cref{sec:hopf-trimodules-lax-mod-funs}.

\begin{proposition}\label{EmperorDaibazaal}
  The diagram
  \[
    \begin{mytikzcd}[ampersand replacement=\&]
      {\mathsf{Str}\cat{V}\mathsf{Mod}(\cat{V}, \cat{V})} \&\& {\mathsf{LexfLax}\cat{V}\mathsf{Mod}(\cat{V}, \cat{V})} \\
      {\cat{V}^{\tensorop}} \&\& \Tetramod[B][B][][B]
      \arrow[hook, from=1-1, to=1-3]
      \arrow["\simeq", from=1-1, to=2-1]
      \arrow["\simeq", from=1-3, to=2-3]
      \arrow[hook, from=2-1, to=2-3]
    \end{mytikzcd}
  \]
  commutes up to a monoidal natural isomorphism.
  Its upper horizontal arrow is the natural inclusion functor,
  its left vertical arrow is the equivalence of \cref{prop:bicategorical-yoneda-correspondence-monads-and-algebras},
  its right vertical arrows is the equivalence of \cref{thm:hopf-trimodules-are-lax-monoidal-functors},
  and its lower horizontal arrow is the functor of \cref{bestoftimes}.
\end{proposition}
\begin{proof}
  Chasing a functor \(F\) in the diagram, we find the following:
  \[
    \begin{mytikzcd}[ampersand replacement=\&]
      F \&\&\& F \\
      {F \Bbbk_{\mathrm{triv}}} \& {B \otimes F \Bbbk_{\mathrm{triv}}} \& {F(B \otimes \Bbbk_{\mathrm{triv}})} \& {FB}
      \arrow[maps to, from=1-1, to=1-4]
      \arrow[maps to, from=1-1, to=2-1]
      \arrow[maps to, from=1-4, to=2-4]
      \arrow[maps to, from=2-1, to=2-2]
      \arrow["\simeq", from=2-2, to=2-3]
      \arrow["\simeq", from=2-3, to=2-4]
    \end{mytikzcd}
  \]
  It is easy to verify that the isomorphism indicated in the bottom row of this chase is natural in \(F\), and also monoidal.
\end{proof}

\begin{corollary}
  If the functor \(B \otimes \blank\) of \cref{bestoftimes} is an equivalence,
  then \(B\) admits a twisted antipode.
\end{corollary}
\begin{proof}
  By \cref{EmperorDaibazaal}, \(B \otimes \blank\) is an equivalence if and only if the embedding
  \[
    \mathsf{Str}\cat{V}\mathsf{Mod}(\cat{V}, \cat{V}) \longhookrightarrow \mathsf{LexfLax}\cat{V}\mathsf{Mod}(\cat{V}, \cat{V})
  \]
  is an equivalence,
  so assume that it is.

  For \(M \in \tetramodfd[B]\), the right adjoint of \(\blank \kotimes M\) is finitary by \cref{saftlex},
  since \(\blank \kotimes M\) sends finite\hyp{}dimensional comodules to finite\hyp{}dimensional comodules.
  The right adjoint is then a finitary left exact lax \(\cat{V}\)\hyp{}module functor.
  By assumption it is strong, and thus of the form \(\blank \otimes \ld{\!{}M}\), for \(\ld{\!{}M} \in \cat{V}\).
  Since the fibre functor \(U\from \cat{V} \to \kVect\) is strong monoidal,
  we have \(\ld{\!{}M} \cong M^{\ast}\).
  This shows that \(\tetramodfd[B]\) is left rigid;
  by \cref{antipodeleftrigid}, \(B\) admits a twisted antipode.
\end{proof}

\begin{proposition}
  If a bialgebra \(B\) admits a twisted antipode,
  then the functor \(B \otimes \blank\) of \cref{bestoftimes} is an equivalence.
\end{proposition}
\begin{proof}
  Since the functor \(\cat{V} \longhookrightarrow \mathsf{LexfLax}\cat{V}\mathsf{Mod}(\cat{V}, \cat{V})\) is always fully faithful,
  so is the functor \(B \otimes \blank\) by \cref{EmperorDaibazaal}.
  It remains to show that \(B \otimes \blank\) is essentially surjective.

  Following~\cite[Section~1.9.4]{montgomery93:hopf},
  for \(X \in \Tetramod[][B][][B]\)
  the composite map
  \[
    \Gamma\from X
    \xrightarrow{\ \Delta_{X}\ } X \otimes B
    \xrightarrow{\ \mathtt{t}\otimes B\ } X^{\mathrm{co}B} \otimes B
    \xrightarrow{\sigma_{X^{\mathrm{co}B},B}} B \otimes X^{\mathrm{co}B}
  \]
  is an isomorphism in \(\Tetramod[][B][][B]\),
  where \(X^{\mathrm{co}B} = \{\,m \in X \mid \Delta^{r}_{X}(m) = m \otimes 1\,\}\)
  and \(\mathtt{t}\) is the projection diagrammatically represented by
  \[
    \tikzfig{the-projection-diagrammatically-represented-by}
  \]

  The inverse of \(\Gamma\) is given by the restriction \({\nabla_{X}|}_{\mathrm{co}B}\from B \otimes X^{\mathrm{co}B} \to X\) of the left \(B\)\hyp{}action on \(X\).

  For \(X \in \Tetramod[B][][B]\),
  the space \(X^{\mathrm{co}B}\) of coinvariants is a left \(B\)\hyp{}subcomodule of \(X\),
  since \(X\) is a bicomodule.
  It is easy to check that \({\nabla_{X}|}_{\mathrm{co}B}\from B \otimes X^{\mathrm{co}B} \to X\) is a left \(B\)\hyp{}comodule morphism,
  and thus a morphism in \(\Tetramod[B][B][][B]\),
  making the functor \(B \otimes \blank\) essentially surjective.
\end{proof}

\begin{corollary}\label{cor:twisted-anti-iff-bestoftimes}
  A bialgebra \(B\) admits a twisted antipode if and only if the functor \(B \otimes \blank\) of \cref{bestoftimes} is an equivalence.
\end{corollary}

\begin{remark}
  Here we see that it is crucial
  that \cref{sec:hopf-trimodules-lax-mod-funs} identifies Hopf trimodules with \emph{finitary} lax module functors.
  Given an infinite\hyp{}dimensional comodule \(M\) over \(B\),
  the functor \(\mathrm{Hom}_{\Bbbk}(M, \blank)\from \Tetramod[B] \to \Tetramod[B]\)
  is endowed with lax \(\Tetramod[B]\)\hyp{}structure
  by virtue of \cref{thm:comodule-adjunctions-correspond-to-strong-comodule-functors},
  but is not finitary itself.
\end{remark}

\section{Fusion operators for Hopf monads}

\textsc{Similarly to our study of} \cref{HausserNill} in \cref{sec:Hausser-Nill},
the results of \Bruguieres, Lack, and Virelizier in~\cite{Bruguieres2011}
can be used to characterise rigidity of a monoidal category.

\begin{lemma}\label{delusion}
  Let \(\adj{F}{U}{\cat{C}}{\cat{D}}\) be an oplax monoidal adjunction.\marginnote{\textnormal{%
      A variant of this result
      using ``Doi--Koppinen data''
      appears in~\cite[Section~5.5]{kowalzig15:cyclic}.%
    }}
  The strong monoidal structure of \(U\) turns \(\cat{C}\) into a \(\cat{D}\)\hyp{}module category, by defining
  \[
    \blank \lact \bblank \;\defeq\; U(\blank) \otimes \bblank.
  \]

  With respect to this \(\cat{D}\)\hyp{}module structure,
  the bimonad \(T \defeq UF\) on \(\cat{C}\) becomes an oplax \(\cat{D}\)\hyp{}module monad.
  The coherence morphism is given by
  \begin{equation}\label{confusion}
    T_{\mathsf{a};d,c} \from T(d \lact c) = T(Ud \otimes c)
    \xrightarrow{T_{2; Ud,c}} TUd \otimes Tc
    \xrightarrow{U\varepsilon_{d} \otimes Tc} Ud \otimes Tc,
  \end{equation}
  for all \(c \in \cat{C}\) and \(d \in \cat{D}\).
  In particular, for \(x \in \cat{C}\),
  the coherence morphism \(T_{\mathsf{a};Fx,c}\) is precisely the right fusion operator for \(T\) at \((x, c)\).
  In other words, \(T_{\mathsf{a};Fx,c} = T_{\mathsf{rf};x,c}\).
\end{lemma}
\begin{proof}
  The first statement is a well-known result on transport of structure.
  For \(y \in \cat{D}\),
  the \(\cat{D}\)\hyp{}module structure on \(U\) is given by the morphisms
  \[
    U_{\mathsf{m}}\from d \lact Uy = Ud \otimes Uy \iso U(d \otimes y) = U(d \lact y).
  \]

  For the second statement,
  observe that by \cref{porism:doctrinal-module-adjunctions},
  \(F\) inherits an oplax \(\cat{D}\)\hyp{}module functor structure
  from the strong \(\cat{D}\)\hyp{}module structure on \(U\).
  The oplax \(\cat{D}\)\hyp{}module functor on \(T\) is given by
  the upper part of the following diagram,
  where the lower part is included clarity:
  \[
    \begin{mytikzcd}[ampersand replacement=\&]
	{UF(Ud \otimes c)} \& {U(FUd \otimes Fc)} \\
	{T(d \lact c)} \& {U(d \otimes Fc)} \& {Ud \otimes UFc} \\
	{UF(d \lact c)} \& {U(d \lact Fc)} \& {Ud \otimes Tc} \\
	\& {d \lact UFc} \& {d \lact Tc}
	\arrow["{U F_{2;Ud,c}}", from=1-1, to=1-2]
	\arrow[equals, from=1-1, to=2-1]
	\arrow["{U(\varepsilon_{d} \otimes Fc)}", from=1-2, to=2-2]
	\arrow[equals, from=2-1, to=3-1]
	\arrow["{U_{\mathsf{m};d,Fc}}", from=2-2, to=2-3]
	\arrow[equals, from=2-2, to=3-2]
	\arrow[equals, from=2-3, to=3-3]
	\arrow[from=3-1, to=3-2]
	\arrow[from=3-2, to=4-2]
	\arrow[equals, from=3-3, to=4-3]
	\arrow[equals, from=4-2, to=4-3]
      \end{mytikzcd}
    \]

  To see that this morphism equals that defined in \cref{confusion},
  observe that the following diagram of functors
  commutes by naturality of \(U_2\):
  \[
    \begin{mytikzcd}[ampersand replacement=\&]
      {U(FU \otimes F)} \&\& {UFU \otimes UF} \\
      {U(\Id_{\cat{D}} \otimes F)} \&\& {U \otimes UF}
      \arrow["{U_{2;FU,F}}", from=1-1, to=1-3]
      \arrow["{U(\varepsilon \,\otimes\, F)}"', from=1-1, to=2-1]
      \arrow["{U\varepsilon \,\otimes\, UF}", from=1-3, to=2-3]
      \arrow["{U_{2;\Id, F}}"', from=2-1, to=2-3]
    \end{mytikzcd}
  \]

  Since
  by definition we have that \(\mu = U \varepsilon F\),
  substituting \(d = Fx\) in \cref{confusion}
  yields \(T_{\mathsf{a};Fx,c} = T_{\mathsf{rf};x,c}\),
  which proves the result.
\end{proof}

\begin{lemma}\label{sinfusion}
  Let \(\cat{C}\) be a category,
  \(f,g\from x \to y\) a pair of parallel morphisms in \(\cat{C}\),
  such that the coequaliser \(\coeq(f,g)\) exists.
  If \(G,H\from \cat{C} \to \cat{D}\) is a pair of functors preserving that coequaliser,
  \(\alpha\from G \nt H\) a natural transformation,
  then
  \(\alpha_{x}\) and \(\alpha_{y}\) being invertible implies that so is \(\alpha_{\coeq(f,g)}\).
\end{lemma}
\begin{proof}
  Let \(c \defeq \coeq(f,g)\)
  and write \(p\from y \longtwoheadrightarrow c\) for the canonical projection.
  We are in the following situation:
  \[
    \begin{mytikzcd}[ampersand replacement=\&]
      Gx \& Gy \& Gc \\
      Hx \& Hy \& Hc
      \arrow["Gf", shift left=1, from=1-1, to=1-2]
      \arrow["Gg"', shift right=1, from=1-1, to=1-2]
      \arrow["{\alpha_x}"', from=1-1, to=2-1]
      \arrow["{Gp}", two heads, from=1-2, to=1-3]
      \arrow["{\alpha_y}", from=1-2, to=2-2]
      \arrow["Hf", shift left=1, from=2-1, to=2-2]
      \arrow["Hg"', shift right=1, from=2-1, to=2-2]
      \arrow["{Hp}", two heads, from=2-2, to=2-3]
    \end{mytikzcd}
  \]
  Then we calculate
  \[
    Hp \circ \alpha_y \circ Gf
    = Hp \circ Hf \circ \alpha_x
    = Hp \circ Hg \circ \alpha_x
    = Hp \circ Gg \circ \alpha_x,
  \]
  which,
  by the universal property of \(Gc\),
  implies the existence of a unique arrow \({!}\from Gc \to Hc\)
  that is equal to \(\alpha_c\).
  Since \(\alpha_x\) and \(\alpha_y\) are invertible
  we also obtain a unique arrow \(Hc \to Gc\),
  which is easily seen to be inverse to \(\alpha_c\).
\end{proof}

\begin{proposition}\label{diffusion}
  The right fusion operators of bimonad \(T\) on a monoidal category \(\cat{C}\) are invertible
  if and only if
  the coherence morphisms for the oplax \(\cat{C}^T\)\hyp{}module structure on \(T\)
  of \cref{delusion} are invertible.

  In other words, \(T\) is right Hopf
  if and only if
  it is a strong \(\cat{C}^T\)\hyp{}module monad.
\end{proposition}
\begin{proof}
  Since, by \cref{delusion}, the right fusion operators for \(T\) are
  a special case of the coherence morphisms for the \(\cat{C}^T\)\hyp{}module functor structure on \(T\),
  it is immediate that \(T\) being a strong \(\cat{C}^T\)\hyp{}module monad implies invertibility of the right fusion operators.

  To show the converse,
  assume that the right fusion operators are invertible,
  in particular, by \cref{delusion},
  the morphisms \(T_{\mathsf{a};Fx,c}\) are invertible.
  Recall any \(T\)\hyp{}algebra \((x, \alpha) \in \cat{C}^T\)
  is the coequaliser of the diagram
  \[
    \begin{mytikzcd}[ampersand replacement=\&]
      {T^{2}x} \& {Tx} \& x
      \arrow["{\mu_{x}}", shift left=1, from=1-1, to=1-2]
      \arrow["{T \alpha}"', shift right=1, from=1-1, to=1-2]
      \arrow["{\alpha}", from=1-2, to=1-3]
    \end{mytikzcd}
  \]
  in \(\cat{C}^T\).
  In other words, \(x\) is a \(U\)\hyp{}split coequaliser of morphisms from \(FTx\) to \(Fx\),
  since \(Fx\) is the free \(T\)\hyp{}module on \(x\).

  Consider the natural transformation \(T_{\mathsf{a};\blank,c}\from T(\blank \lact c) \to \blank \lact Tc\).
  We have
  \[
    T(\blank \lact c) = T(U(\blank) \otimes c) = (\blank \otimes c) \circ U,
  \]
  which means that \(T(\blank \lact c)\) preserves the coequaliser \(\coeq(T\alpha, \mu_{x})\)
  by first sending it to a split coequaliser under \(U\),
  and then preserving it, since split coequalisers are absolute.

  In a similar fashion, \(\blank \lact Tc = U(\blank) \otimes Tc = (\blank \otimes Tc) \circ U\) also preserves this coequaliser.
  The invertibility of \(T_{\mathsf{a};x,c}\) now follows from
  the assumed invertibility of
  \(T_{\mathsf{rf};x,c} = T_{\mathsf{a};Fx,c}\)
  and \(T_{\mathsf{rf}; Tx, c} = T_{\mathsf{a};FTx,c}\),
  by \cref{sinfusion}.
\end{proof}

%%% Local Variables:
%%% mode: LaTeX
%%% TeX-master: "../main"
%%% TeX-command-extra-options: "-shell-escape -fmt=prec.fmt -file-line-error -synctex=1"
%%% End:
